\section*{Adaptive single-use payment for Boolean root KL rows}

Let $f:\{-1,1\}^n\to\{-1,1\}$ and let $\mathcal T$ be any physical
restriction tree.  At an internal node $v$, the event $E_v$ is the
cylinder determined by the preceding answers, $p_v=\Pr(E_v)$, and the
tree queries a coordinate $i_v$ not fixed on $E_v$.  Write $f_v$ for the
restricted Boolean function and
\[
 d_v=\widehat {f_v}(\{i_v\})
     =\mathbb E[f(x)x_{i_v}\mid E_v].
\]
No physical coordinate is queried twice on a path, but the next query may
depend arbitrarily on all preceding answers.  The tree may stop early.

\paragraph{Theorem (adaptive root nonreuse).}
If $K_{i_v}^{(0)}(f_v)$ denotes the empty-upper-set row of the exact local
identity
\[
 \Delta_{i_v}U(f_v)=B_{i_v}(f_v)-\Gamma_{i_v}(f_v)
                     +K_{i_v}(f_v),
\]
then, with $c_0=2/\ln2$,
\[
 \boxed{\sum_{v\in\mathcal T^\circ}p_vK_{i_v}^{(0)}(f_v)
 \le c_0\sum_{v\in\mathcal T^\circ}p_vd_v^2
 \le c_0\operatorname {Var}(f)\le c_0 I(f).}       \tag{1}
\]
If the tree reveals every coordinate on every path, the middle inequality
is equality.  Thus the sharp root-row payment is single-use even under
fully adaptive physical restriction; it has no depth, arity, overlap,
bias, or minimum-gap loss.

\paragraph{Proof of the pointwise payment.}
Boolean convolution for $f_v^2=1$ gives zero signed upper correlation in
the empty row.  Put $a_v=\mathbb E[f\mid E_v]$ and
$m_v=a_v^2+d_v^2$.  The sharp scalar inequality
\[
1-h_2\left({1+x\over2}\right)
\le {1-\sqrt{1-x^2}\over\ln2},\qquad |x|\le1,
\]
therefore gives
\[
K_{i_v}^{(0)}(f_v)
=m_v\left[1-h_2\left({1+2a_vd_v/m_v\over2}\right)\right]
\le {2\min(a_v^2,d_v^2)\over\ln2}
\le c_0d_v^2.                                      \tag{2}
\]
This uses genuine Booleanity at every physical child, not an arbitrary
spectral law.

\paragraph{Orthogonal single-use ledger.}
For each internal node define
\[
 \psi_v(x)={\mathbf1_{E_v}(x)x_{i_v}\over\sqrt{p_v}}.
\]
These functions have norm one.  Two incomparable nodes have disjoint
supports.  If $w$ is a strict descendant of $v$, then $x_{i_v}$ is fixed
on $E_w$, whereas the still-unqueried $x_{i_w}$ is fair conditionally on
$E_w$; hence $\langle\psi_v,\psi_w\rangle=0$.  Thus the $\psi_v$ are
orthonormal and
\[
 \langle f,\psi_v\rangle=\sqrt{p_v}\,d_v.
\]
They are also orthogonal to the constant function.  Bessel's inequality
gives
\[
 \sum_vp_vd_v^2\le\|f-\mathbb Ef\|_2^2=\operatorname {Var}(f).
\]
A complete tree supplies the usual complete Boolean Haar basis, giving
equality.  Combining this with (2) proves (1).

\paragraph{Exact implication for the GU programme.}
At one function the already proved sharp bound is
\[
 \sum_iK_i^{(0)}\le c_0I.
\]
Consequently, for any $a\ge c_0$ and $0\le b<1$, the still-unpaid global
inequality
\[
 \sum_iK_i^{(+)}
 \le(a-c_0)I+bU+\sum_i(\Gamma_i-B_i)                \tag{3}
\]
on tensor-prime Boolean functions implies $G\le aI+bU$ universally and
therefore
\[
 H\le\left(2+{2a\over1-b}\right)I.
\]
The new content of (1) is that recursive or stopping-tree attempts to
prove (3) may transport the paid empty rows through adaptive physical
children without allocating another $c_0I$ at every depth: the entire
selected-coordinate root ledger costs at most $c_0\operatorname{Var}(f)$.

Equivalently, a stopped-potential proof now needs to pay only the
nonempty-row transport error and its terminal $U$ boundary.  It must not
identify a parent nonempty row with a descendant root row without proving
the required comparison; that comparison is the first remaining line.
No universal pair $(a,b)$ is established here, and full FEI remains open.

